\documentclass[11pt,a4paper]{article}
\usepackage[italian]{babel}
\usepackage[margin=2.5cm]{geometry}
\usepackage{parskip}
\usepackage{amsmath, amssymb, amsthm}
\usepackage[most]{tcolorbox}
\usepackage{array}
\usepackage{tikz}
\usetikzlibrary{decorations.pathreplacing}
\usepackage{hyperref}

% --- Riquadri con bordo nero, senza numero ---
% Uso: \begin{definizione} ... \end{definizione}
%      \begin{teorema}[Nome] ... \end{teorema}  (il nome è facoltativo)
\tcbset{nota/.style={colback=white, colframe=black, boxrule=0.5pt,
  sharp corners}}
\NewTColorBox{definizione}{}{nota,
  before upper={\textbf{Definizione.}\ }}
\NewTColorBox{teorema}{o}{nota,
  before upper={\textbf{Teorema\IfValueT{#1}{ (#1)}.}\ }}
\NewTColorBox{proposizione}{o}{nota,
  before upper={\textbf{Proposizione\IfValueT{#1}{ (#1)}.}\ }}
\NewTColorBox{assioma}{o}{nota, boxrule=1pt,
  before upper={\textbf{Assioma\IfValueT{#1}{ (#1)}.}\ }}

% --- Ambienti semplici (senza riquadro) ---
\theoremstyle{definition}
\newtheorem*{esempio}{Esempio}
\newtheorem*{osservazione}{Osservazione}
\newtheorem*{esercizio}{Esercizio}

% V e F nelle tavole di verità
\newcommand{\V}{\textsf{V}}
\newcommand{\F}{\textsf{F}}

\title{Operazioni tra funzioni}
\author{Marco Cerbella}
\date{Lezione 4 -- 5 ottobre 2026}

\begin{document}
\maketitle

\section{Operazioni aritmetiche tra funzioni}

Date due funzioni $f : D(f) \subseteq \mathbb{R} \to \mathbb{R}$ e $g : D(g) \subseteq \mathbb{R} \to \mathbb{R}$ si definiscono:
\begin{itemize}
    \item \emph{somma}: $(f + g)(x) = f(x) + g(x)$;
    \item \emph{differenza}: $(f - g)(x) = f(x) - g(x)$;
    \item \emph{prodotto}: $(fg)(x) = f(x)\,g(x)$;
    \item \emph{quoziente}: $\left(\dfrac{f}{g}\right)(x) = \dfrac{f(x)}{g(x)}$.
\end{itemize}
Il dominio delle prime tre funzioni è $D(f) \cap D(g)$. Il dominio del quoziente è
\[
D(f) \cap \{x \in D(g) : g(x) \neq 0\}.
\]

\section{Funzione massimo e funzione minimo}

Date $f : D(f) \subseteq \mathbb{R} \to \mathbb{R}$ e $g : D(g) \subseteq \mathbb{R} \to \mathbb{R}$ si definiscono:
\begin{itemize}
    \item \emph{massimo}: $\max\{f, g\}(x) = \max\{f(x), g(x)\}$;
    \item \emph{minimo}: $\min\{f, g\}(x) = \min\{f(x), g(x)\}$.
\end{itemize}
Il dominio di queste funzioni è $D(f) \cap D(g)$. Nel grafico, $\max\{f, g\}$ segue in ogni punto la più alta tra le due curve, $\min\{f, g\}$ la più bassa.

\end{document}
